Given an ambiguous financial question instance (example in Table~\ref{tab:example}), it should contain the ambiguous question over a company filing $Q$, the disambiguating context $C$ which is the minimal set of scope or definitional constraints that collapses $Q$ into the unique intended answer $A$ under intended interpretation $\mathcal{I}_{\text{int}}$. The ambiguous question typically contains correct answer $A_{d}\neq A$ under the default interpretation $\mathcal{I}_{\text{def}}$ which a non-expert would assume when reading $Q$. Specifically, $\mathcal{I}_{\text{int}},\mathcal{I}_{\text{def}}$ are structured records specifying the financial ambiguity taxonomy defined below. 
\begin{table}[h]
\centering
\caption{This example of ambiguous financial question contains three plausible entity scopes, i.e. consolidated total and two reportable segments, leading to an entropy of $1.58$ bits. Without the disambiguating context $C$, a model's default interpretation leads to the wrong answer.}
\label{tab:example}
\small
\begin{tabular}{p{2.7cm}p{4.7cm}}
\toprule
Field & Value \\
\midrule
$Q$ &
  What was Meta Platforms' operating income? \\
$C$ &
  Total company (consolidated) results under U.S.\ GAAP for the year
  ended December 31, 2023. \\
$A$ (intended) &
  \$46.75B \emph{[consolidated GAAP]} \\
$A_d$ (default) &
  \$62.87B \emph{[Family of Apps segment]} \\
\midrule
Ambiguity category & Entity scope \\
$H_0$ & 1.58 bits \\
\midrule
Intended span &
  ``Income from operations for 2023 was \$46.75 billion\ldots'' \\
Default span &
  ``Family of Apps\ldots Income from operations \$62,871\ldots'' \\
\bottomrule
\end{tabular}
\end{table}

\textbf{Financial ambiguity taxonomy.}
\label{sec:taxonomy}
We derive five categories of financial ambiguity from standard
accounting and disclosure practice (Table~\ref{tab:axes}). Each instance should normally exercise one primary category with one or two optional secondary categories in case of additional ambiguity. 

\begin{table}[h]
\centering
\caption{Five-category financial ambiguity taxonomy.}
\label{tab:axes}
\small
\begin{tabular}{p{2.0cm}p{5.5cm}}
\toprule
Category & Definition \\
\midrule
\textbf{Temporal scope} &
  Ambiguity over which time period is intended: fiscal year ending
  September vs. calendar year; Q3 results vs. full year; TTM vs.
  point-in-time. \\
\addlinespace
\textbf{Metric definition} &
  Ambiguity over accounting basis: GAAP net income vs. non-GAAP adjusted
  income; organic revenue vs. as-reported; operating EBITDA vs. net income. \\
\addlinespace
\textbf{Entity scope} &
  Ambiguity over consolidation level: segment vs. consolidated total;
  parent-attributable vs. full consolidated; Class A vs. Class B shares;
  A-shares vs. H-shares. \\
\addlinespace
\textbf{Filing vintage} &
  Ambiguity over filing version: original 10-K vs. amended 10-K/A;
  as-reported vs. restated; preliminary earnings release vs. final filing. \\
\addlinespace
\textbf{Recognition policy} &
  Ambiguity over accounting policy: revenue recognized point-in-time
  vs. over time; gross vs. net revenue; capitalized vs. expensed R\&D. \\
\bottomrule
\end{tabular}
\end{table}

\textbf{Disambiguation entropy.} We quantify how ambiguous an instance is using the prior confidence entropy $H_{0} = \log_{2} k$, where $k$ denotes the number of plausible interpretations of $Q$ without clarifying interaction, e.g. the example in Table~\ref{tab:example} has three plausible interpretations, hence contributing to an entropy of $H_0 = \log_2 3 \approx 1.58$ bits. This further serves as our metric to evaluate difficulty and efficiency in Sec. \ref{sec:metrics}. While both $A$ and $A_{d}$ are defensible from publicly available filings, an agent should be able to identify the ambiguity by reading $Q$ alone without $C$, where this pairing between default and intended interpretation generalizes the target-distractor design from \interactcomp~\cite{interactcomp2026} to the financial domain. 

\textbf{Agent Interaction Model.} We follow \interactcomp's ReAct framework~\citep{yao2023react} where the agent may issue three types of actions, namely \textbf{search}($Q$) to retrieve a passage from the filing corpus, \textbf{interact}($Q$) to present a yes-or-no question to the simulated user who answers based on $C$, and \textbf{answer}($Q$) to submit a final answer. Specifically, we evaluate agents in three primary modes (Table~\ref{tab:modes}) and four
ablation modes. We additionally include a category-conditioned interaction mode, described in
Appendix~\ref{sec:axisreact}.

\begin{table}[h]
\centering
\caption{Evaluation modes.}
\label{tab:modes}
\small
\begin{tabular}{lll}
\toprule
Mode & Actions & Description \\
\midrule
Answer-only & answer & Pure parametric recall \\
Answer+Search & search, answer & Oracle retrieval augmented \\
Answer+Search+Interact & search, interact, answer & Standard full interaction \\
\midrule
Always-ask & \multicolumn{2}{l}{Must ask $\geq 1$ question before answering} \\
Category-oracle & \multicolumn{2}{l}{Told the primary category, not the value} \\
Template-oracle & \multicolumn{2}{l}{Fixed human-written question per category} \\
Enumerate & \multicolumn{2}{l}{Lists all interpretations and answers} \\
\bottomrule
\end{tabular}
\end{table}